\chapter{Экспериментальная оценка, интеграция и анализ эффективности модуля ранжирования кандидатов}
\label{chapter:eval}
\section{Сравнение подхода на основе моделей с механизмом самовнимания с классическими системами ранжирования}

Перед интерпретацией экспериментальных метрик необходимо явно определить, с какими классами классических систем ранжирования сопоставляется выбранный в работе подход на базе двунаправленного кодировщика представлений (Bidirectional Encoder Representations from Transformers), модели векторных представлений предложений (Sentence Bidirectional Encoder Representations from Transformers) и архитектуры двойного кодировщика (dual encoder architecture). В задаче автоматизированного подбора персонала запросом фактически является текст вакансии, а документами - резюме кандидатов, поэтому исходная постановка близка к классическому информационному поиску. В системах автоматизированного подбора персонала исторически применялись векторная модель поиска, лексические схемы взвешивания терминов, вероятностные модели релевантности, например вероятностная функция ранжирования документов под названием BM25, ссылочные или графовые методы, а также методы обучения ранжированию над заранее рассчитанными признаками \cite{src08, src11, src25}.

Векторная модель поиска представляет документы и запросы как векторы терминов, после чего упорядочивает документы по мере близости, например по косинусному сходству \cite{src25}. Модель частоты термина с обратной частотой документа (Term Frequency-Inverse Document Frequency, TF-IDF) конкретизирует эту идею через повышение веса редких и потенциально более специфичных терминов, что делает метод прозрачным и удобным контрольным ориентиром для текстового поиска \cite{src26}. В сфере подбора персонала такая логика хорошо работает для точных совпадений названий технологий, должностей и сертификатов, однако слабо реагирует на синонимию и вариативность профессионального языка: резюме и вакансия могут описывать один и тот же опыт разными словами. Вероятностная функция ранжирования документов под названием BM25, развивая вероятностную парадигму информационного поиска, добавляет насыщение частоты термина и нормализацию по длине документа, поэтому является более устойчивой классической функцией ранжирования для коллекций неоднородных текстов \cite{src11}. Тем не менее функция под названием BM25 остается в основном лексическим методом: она оценивает совпадение наблюдаемых терминов, а не контекстуальную близость смыслов.

Подход на основе двунаправленного кодировщика (Bidirectional Encoder Representations from Transformers) отличается тем, что строит контекстные плотные представления текста, а вариант с моделью векторных представлений предложений (Sentence Embeddings using Siamese Bidirectional Encoder Representations from Transformers Networks) позволяет сравнивать вакансии и резюме как семантические векторы, полученные независимым преобразованием текстов \cite{src05, src06}. Для данной работы это существенно, поскольку тексты резюме и вакансий в системах автоматизированного подбора персонала часто содержат неполные формулировки, разные уровни детализации и разные способы описания одних и тех же навыков. Выбранная архитектура двойного кодировщика (dual encoder architecture) на базе \path|cointegrated/rubert-tiny2| не заменяет все классические методы, но решает именно ту часть задачи, где простое совпадение терминов недостаточно: она формирует числовую меру близости текстовой пары текст вакансия и текст резюме, и может пройти доменную тонкую настройку на исторических событиях подбора персонала целевого корпуса размеченных данных.

Графовые методы, например PageRank, имеют иной источник сигнала: они ранжируют объекты по структуре ссылок или связей между ними \cite{src27}. Для веб-поиска или научного цитирования это естественно, но в рассматриваемой задаче резюме и вакансии не образуют полноценного ссылочного графа, сопоставимого с графом веб-страниц. Поэтому PageRank не является основным контрольным методом для текущего модуля семантического сопоставления. В перспективе графовые признаки могли бы использоваться как вспомогательные сигналы при наличии устойчивого графа взаимодействий кандидатов, вакансий, работодателей и откликов, однако такая постановка задачи выходит за пределы реализованного прототипа модуля ранжирования цифровых профилей соискателей.

Классические методы обучения ранжированию (Learning to Rank), включая метод попарного обучения ранжированию под названием RankNet, метод градиентного бустинга деревьев решений для ранжирования под названием LambdaMART и другие варианты градиентного бустинга для ранжирования, решают более общую задачу: они обучают итоговую ранжирующую функцию по наборам признаков и разметке предпочтений или релевантности \cite{src08, src09, src28}. Их преимущество состоит в возможности объединять разнородные сигналы: оценку вероятностной функции ранжирования документов под названием BM25, модель частоты термина с обратной частотой документа (Term Frequency-Inverse Document Frequency, TF-IDF), косинусное сходство (cosine similarity), структурные совпадения навыков, стаж, зарплатные ожидания, историю откликов и поведенческие метрики. Ограничение для текущей работы заключается в том, что полноценная модель обучения ранжированию (Learning to Rank) требует специально сформированных списков кандидатов на каждую вакансию, достаточного числа градуированных меток и отдельной процедуры оптимизации под метрики ранжирования. Поэтому переранжирование с обучением ранжированию (Learning to Rank) в тексте рассматривается как логичное развитие второго этапа, тогда как реализованная часть работы ограничивается архитектурой двойного кодировщика (dual encoder architecture) и сортировкой по косинусному сходству.

С учетом этого различия описанный подход на базе моделей с механизмом самовнимания в данной работе следует рассматривать не как универсальную замену всем классическим методам, а как способ построения более содержательного признака релевантности для текстовой пары текст вакансии и текст резюме. Лексические функции, такие как модель частоты термина с обратной частотой документа (Term Frequency-Inverse Document Frequency, TF-IDF) и вероятностная функция ранжирования документов под названием BM25, остаются полезными как дешевые и интерпретируемые контрольные методы: они позволяют быстро проверить, насколько далеко продвигает качество простое совпадение терминов, и могут применяться в гибридных схемах первичной фильтрации. Однако именно в текстах подбора персонала их ограничение проявляется особенно заметно, поскольку релевантность часто выражается не буквальным совпадением слов, а согласованностью содержания вакансии и резюме: совпадением названия должности и уровня позиции, пересечением требуемых и указанных в резюме навыков и технологий, а также сопоставимостью описанных задач, обязанностей и достижений. Кандидат может быть релевантен по смыслу, даже если в текстах использованы разные формулировки. Ссылочные методы наподобие PageRank в данной постановке имеют меньший вес: они естественны для веб-графов и сетей цитирования, но не описывают напрямую смысловое соответствие вакансии и резюме. Методы обучения ранжированию (Learning to Rank), напротив, являются содержательным направлением дальнейшего развития: при наличии достаточного числа списков кандидатов и градуированных оценок они могут объединить оценку вероятностной функции ранжирования документов под названием BM25, семантический показатель, полученный архитектурой двойного кодировщика (dual encoder architecture), структурные признаки профиля и поведенческие сигналы системы автоматизированного подбора персонала в единую ранжирующую функцию. Поэтому выбор двунаправленного кодировщика (Bidirectional Encoder Representations from Transformers) и модели векторных представлений предложений на базе сиамских кодировочных сетей (Sentence Embeddings using Siamese Bidirectional Encoder Representations from Transformers Networks) в текущей реализации обусловлен не отказом от классического информационного поиска, а тем, что в рамках настоящей работы на первом этапе поставлена задача получить и экспериментально проверить воспроизводимое семантическое ранжирование для русскоязычных текстовых текст вакансия и текст резюме. Построение полноценной гибридной модели ранжирования с обучением ранжированию (Learning to Rank) вынесено за пределы настоящей работы и предполагает использование этого показателя как одного из входных признаков.

\section{Проведение экспериментальной оценки качества алгоритма ранжирования по нормализованной дисконтированной кумулятивной выгоде, точности, полноте и средней обратной позиции}

Модуль оценивался в формате офлайн-валидации на пересобранном наборе данных, сформированном из первоисточников агрегатора вакансий и резюме. После очистки и пересборки связанных итоговая выборка составила 5 218 примеров, из которых 4 176 использовались для обучения, 521 для валидации и 521 для финального тестирования. Основной сравнительный протокол включает ровно три модели: вероятностную функцию ранжирования документов под названием BM25 как лексический контрольный метод без обучения; перекрестную модель совместного кодирования текстовой пары под названием \path|cointegrated/rubert-tiny2|, прошедшую тонкую настройку (fine-tuning) на train.tsv и валидацию на val.tsv; модель раздельного кодирования вакансии и резюме (bi-encoder), сохраненную в каталоге \path|experiments/models/bi_encoder_rubert_tiny2| и рассматриваемую как основной предложенный в работе подход.

На уровне отдельных текстовых текст вакансии и текст резюме качество оценивается по средней точности по всей кривой точность-полнота (Average Precision), гармонической средней (F1-score), доле верных решений (Accuracy) и коэффициенту корреляции Мэтьюса (Matthews Correlation Coefficient, MCC), а финальные значения публикуются только на файле test.tsv. Для вероятностной функции ранжирования документов под названием BM25 числовая оценка задается классической формулой релевантности между вакансией как запросом и резюме как документом после регулярной токенизации по словам и приведения текста к нижнему регистру. Для перекрестного кодировщика входом является пара текстов [vacancy\_text, resume\_text], а числовая оценка равна вероятности класса label = 1. Для модели раздельного кодирования (bi-encoder) вакансия и резюме преобразуются в отдельные векторы смысла, а оценка рассчитывается как косинусное сходство (cosine similarity) между этими векторными представлениями. В имеющемся прогоне основной предложенный подход это модель раздельного кодирования текстов вакансии и резюме, после тонкой настройки (fine-tuning) достиг на валидационной выборке данных значения средней точности (Average Precision) 0.9870, лучшей гармонической средней (F1-score) 0.9582, лучшей доли верных решений (Accuracy) 0.9712 и коэффициента корреляции Мэтьюса (MCC) 0.9365; на test модель показала среднюю точность (Average Precision) 0.9921, лучшую гармоническую среднюю (F1-score) 0.9636, лучшую долю верных решений (Accuracy) 0.9750 и коэффициент корреляции Мэтьюса (MCC) 0.9447.

Дополнительно была выполнена оценка метрик ранжирования на полном файле тестовой выборки \path|data/splits/test.tsv| с группировкой по идентификатору вакансии vacancy\_id, что приближает эксперимент к реальному сценарию сортировки кандидатов внутри вакансии. В тестовом срезе содержится 521 пара, 423 уникальные вакансии, 510 уникальных резюме, 177 пар, где кандидат был принят к рассмотрению или приглашен на следующий этап, и 344 пары, где по отклику был зафиксирован отказ. При этом тестовый набор является разреженным: для 304 вакансий в тестовой разметке нет ни одного кандидата, по которому историческое решение специалиста соответствует продвижению на следующий этап, только 59 вакансий содержат минимум двух кандидатов, а нормализованная дисконтированная кумулятивная выгода (Normalized Discounted Cumulative Gain, NDCG) математически определена только для 33 вакансий, где одновременно есть хотя бы два кандидата и хотя бы один исторически продвинутый кандидат. Для вакансий с хотя бы одним таким кандидатом были получены следующие значения: нормализованная дисконтированная кумулятивная выгода на первых трех позициях (Normalized Discounted Cumulative Gain at 3, NDCG@3) 1.0000, точность на первых трех позициях (Precision at 3, Precision@3) 0.9958, полнота на первых трех позициях (Recall at 3, Recall@3) 0.9824 и средняя обратная позиция первого исторически продвинутого кандидата (Mean Reciprocal Rank, MRR) 1.0000. Эти показатели указывают на то, что на сформированном тестовом наборе дообученная модель в большинстве случаев помещает кандидатов с исторически подтвержденным продвижением в верхнюю часть выдачи.

Таким образом, экспериментальная логика работы состоит не в переборе близких вариантов одной модели, а в сравнении трех разных компромиссов. Вероятностная функция ранжирования документов под названием BM25 показывает силу классического поиска по совпадению слов и терминов; перекрестная модель совместного кодирования текстовой пары (cross-encoder) показывает потенциальную верхнюю границу качества при дорогом совместном анализе каждой пары. Модель раздельного кодирования (bi-encoder) проверяет пригодность масштабируемой семантической архитектуры для интеграции в систему автоматизированного подбора персонала. Полученные результаты следует интерпретировать как офлайн-валидацию на исторических данных, а не как онлайн-эксперимент в рабочем контуре системы автоматизированного подбора персонала под названием Reqcore.

Ограничения экспериментальной части также должны быть зафиксированы явно. Сравнение вероятностной функции ранжирования документов под названием BM25, перекрестной модели совместного кодирования текстовой пары (cross-encoder) и модели раздельного кодирования (bi-encoder) не заменяет валидацию объяснимости специалистами по подбору персонала, аудит предвзятости модели по социальным признакам и нагрузочное тестирование модуля ранжирования на больших наборах цифровых профилей кандидатов. Кроме того, экспериментальная ветка обработки biencoder в текущем коде возвращает только числовую оценку и не формирует объяснения на уровне отдельных критериев, подтверждающие фрагменты (evidence), сильные стороны (strengths) и пробелы соответствия (gaps), поэтому результаты корректно трактовать именно как проверку качества числового сопоставления и ранжирования, а не как оценку полноты объяснимости решения.

\section{Выполнение интеграции созданного модуля в целевую систему автоматизированного подбора персонала}

Интеграция реализована в виде отдельного сервиса программного интерфейса приложения (Application Programming Interface, API), работающего поверх протокола передачи гипертекста (Hypertext Transfer Protocol, HTTP), на фреймворке языка Python под названием FastAPI. Этот сервис взаимодействует с системой автоматизированного подбора персонала под названием Reqcore по протоколу передачи гипертекста (HTTP). В текущем коде сервис на языке Python \path|experiments/scripts/11_api_server.py| предоставляет три маршрута, то есть три адреса программного интерфейса: POST \path|/index/candidate| загружает кандидатские профили во временный набор данных в оперативной памяти, POST \path|/search/rank| ранжирует список кандидатов по оценке семантического соответствия вакансии и резюме, а POST \path|/score| рассчитывает числовую оценку одной пары текст вакансии и текст резюме. Таким образом, фактически реализованная интеграция представляет собой внешний сервис семантического вычисления соответствия резюме требованиям вакансии с хранением профилей в оперативной памяти, а не связку с промышленной векторной базой данных.

Работоспособность контура программного интерфейса подтверждена отдельным сквозным тестом \path|experiments/scripts/16_e2e_integration_test.py|. В этом сценарии через тестовый клиент фреймворка языка Python под названием FastAPI (FastAPI TestClient) последовательно проверяются загрузка кандидатских профилей через POST \path|/index/candidate|, ранжирование через POST \path|/search/rank| и одиночный расчет числовой близости резюме требованиям вакансии через POST \path|/score|. Тест валидирует ответы по протоколу передачи гипертекста (HTTP), контракт в нотации объектов JavaScript (JSON) и наличие телеметрических заголовков X-Processing-Time-Ms и X-Model-Version. Следовательно, на уровне тестового контура подтверждено, что сервис корректно обрабатывает запросы на получение данных, ранжирование и расчет числовой близости для одной текстовой пары в рамках одного программного интерфейса.

Дополнительно была выполнена локальная проверка интеграции с системой автоматизированного подбора персонала под названием Reqcore. Для этого был поднят локальный демонстрационный стенд системы автоматизированного подбора персонала под названием Reqcore через команду Docker Compose, сервис на языке Python для раздельного кодирования вакансии и резюме (bi-encoder) был запущен отдельно с указанием переменной окружения MODEL\_PATH, а в переменной окружения BIENCODER\_URL для серверного процесса Reqcore был задан адрес сервиса ATS Semantic Ranking Engine API, реализованного в файле \path|experiments/scripts/11_api_server.py|. В результате ветка обработки provider === 'biencoder' в файле \path|ats/reqcore/server/api/applications/[id]/analyze.post.ts| была пройдена до конца: система автоматизированного подбора персонала под названием Reqcore вызвала маршрут POST \path|/score|, получила числовое значение match\_score, преобразовала его в значение поля \path|application.score| и сохранила итоговый результат анализа в analysis\_run с провайдером provider = biencoder и моделью model = \path|cointegrated/rubert-tiny2|. Поскольку в этом сценарии не отправляется запрос к большой языковой модели и не генерируется текстовый ответ, служебные поля учета знаков остаются нулевыми: promptTokens = 0 и completionTokens = 0. Тем самым подтверждена не только автономная работа сервиса на фреймворке языка Python под названием FastAPI, но и его фактическое подключение к демонстрационному контуру системы автоматизированного подбора персонала под названием Reqcore.

Важно, что текущая интеграция с системой автоматизированного подбора персонала под названием Reqcore не заменяет весь контур большой языковой модели целиком, а добавляет альтернативный путь только для расчета итоговой числовой близости резюме требованиям вакансии. В стандартном сценарии оценки через большую языковую модель серверный маршрут Reqcore POST \path|/api/applications/[id]/analyze| продолжает вызывать функцию scoreApplication(), обращаться к выбранному провайдеру большой языковой модели и сохранять детализацию оценки по критериям в таблице criterion\_score. Для biencoder-сценария тот же маршрут использует отдельную ветку с вызовом адреса из переменной окружения BIENCODER\_URL и серверного маршрута \path|/score|, после чего в системе автоматизированного подбора персонала под названием Reqcore записывается только итоговая числовая оценка без полей criterion\_score, evidence, strengths и gaps. Следовательно, интеграция корректно описывается как подключение альтернативного модуля семантического ранжирования текстов вакансий и резюме, а не как внедрение полноценной векторной базы данных или полная замена всей логики анализа резюме с использованием большой языковой модели.

\section{Экономическое обоснование внедрения модуля ранжирования}

Экономическая оценка не опирается на логи траты знаков по подписке у провайдеров и не утверждает, что фактический расход или замер рабочего времени специалистов по подбору персонала извлечены из работающей в реальном времени промышленной системы автоматизированного подбора персонала. В коде системы автоматизированного подбора персонала под названием Reqcore предусмотрены поля числа знаков запроса и ответа (promptTokens, completionTokens), а также настраиваемые ценовые параметры (inputPrice, outputPrice); их определения приведены ниже. Пользовательский интерфейс рассчитывает стоимость по формуле C\_ui, также приведённой ниже. В рамках работы эти значения из базы данных системы не выгружались, поэтому ниже приводится расчетный сценарий на публичной форме стоимости запроса к большой языковой модели, измеренном времени вывода модели и явных допущениях для возможности оценки экономики с допустимой долей погрешности.

\begin{equation}
\label{eq:prompttokens-1}
\text{promptTokens} = \text{analysis\_run.prompt\_tokens}
\end{equation}

\begin{equation}
\label{eq:completiontokens-2}
\text{completionTokens} = \text{analysis\_run.completion\_tokens}
\end{equation}

\begin{equation}
\label{eq:inputprice-3}
\text{inputPrice} = \text{input\_price\_per\_1m}
\end{equation}

\begin{equation}
\label{eq:outputprice-4}
\text{outputPrice} = \text{output\_price\_per\_1m}
\end{equation}

\begin{equation}
\label{eq:c_ui-5}
C_{\text{ui}} = \frac{promptTokens}{1{,}000{,}000} \cdot inputPrice \\
         + \frac{completionTokens}{1{,}000{,}000} \cdot outputPrice
\end{equation}

Формула стоимости одного запроса к большой языковой модели описывает расчет соответствия между требованиями конкретной вакансии и текстом резюме одного кандидата:

\begin{equation}
\label{eq:c_llm_pair-6}
C_{\text{llm\_pair}} = (T_{\text{in}} / 1{,}000{,}000) \cdot P_{\text{in}} + (T_{\text{out}} / 1{,}000{,}000) \cdot P_{\text{out}}
\end{equation}

где T\_in и T\_out - число входных и выходных токенов, а P\_in и P\_out - цены за 1 млн входных и выходных токенов соответственно.

Для сценария предварительного топ-K ранжирования через модель раздельного кодирования вакансии и резюме (bi-encoder) экономия внешних вызовов большой языковой модели описывается так:

\begin{equation}
\label{eq:c_full_llm-7}
C_{\text{full\_llm}}(N) = N \cdot C_{\text{llm\_pair}}
\end{equation}

\begin{equation}
\label{eq:c_bi_infra-8}
C_{\text{bi\_infra}}(N) = (L_{\text{bi}}(N) / 3600) \cdot C_{\text{cpu\_hour}}
\end{equation}

\begin{equation}
\label{eq:c_topk_pipeline-9}
C_{\text{topK\_pipeline}}(N, K) = C_{\text{bi\_infra}}(N) + K \cdot C_{\text{llm\_pair}}
\end{equation}

\begin{equation}
\label{eq:savings_api-10}
Savings_{\text{api}}(N, K) = 1 - C_{\text{topK\_pipeline}}(N, K) / C_{\text{full\_llm}}(N)
\end{equation}

\begin{equation}
\label{eq:savings_approx-11}
Savings_{\text{approx}}(N, K) = 1 - K / N
\end{equation}

где $N$ - число кандидатов в исходном наборе данных, $K$ - число кандидатов, передаваемых на дальнейший анализ большой языковой моделью после предварительного ранжирования, L\_bi(N) - задержка локального расчета близости моделью раздельного кодирования для набора из $N$ кандидатов, а C\_cpu\_hour - стоимость часа инстанса с центральным модулем расчётов (CPU). Если инфраструктурная стоимость локального расчета оценки мала относительно вызовов большой языковой модели, то для N = 100 и K = 10 по формуле Savings\_approx получаем 1 - K / N = 90\%.

Задержка полного оценивания с использованием большой языковой модели и предварительного ранжирования моделью раздельного кодирования задается следующими формулами:

\begin{equation}
\label{eq:l_full_llm_seq-12}
L_{\text{full\_llm\_seq}}(N) = N \cdot L_{\text{llm\_pair}}
\end{equation}

\begin{equation}
\label{eq:l_full_llm_parallel-13}
L_{\text{full\_llm\_parallel}}(N) \approx  \operatorname{ceil}(N / P_{\text{parallel}}) \cdot L_{\text{llm\_pair}}
\end{equation}

\begin{equation}
\label{eq:l_bi-14}
L_{\text{bi}}(N) = L_{\text{filter}}(N) + L_{\text{encode\_vacancy}} + L_{\text{encode\_candidates}}(N) + L_{\text{cosine\_topK}}(N)
\end{equation}

\begin{equation}
\label{eq:l_bi_precomputed-15}
L_{\text{bi\_precomputed}}(N) = L_{\text{filter}}(N) + L_{\text{encode\_vacancy}} + L_{\text{cosine\_topK}}(N)
\end{equation}

где L\_llm\_pair - задержка одного запроса к большой языковой модели для пары текст вакансии и текст резюме, P\_parallel - допустимое число параллельных запросов к большой языковой модели, L\_filter(N) - время применения жестких фильтров, L\_encode\_vacancy - время преобразования текста вакансии в векторное представление, L\_encode\_candidates(N) - время преобразования текстов кандидатов в векторные представления, L\_cosine\_topK(N) - время расчета косинусных сходств и выбора топ-K. Оценивание с использованием большой языковой модели дает критериальное объяснимое оценивание, но зависит от внешнего программного интерфейса и переменной стоимости вызова. Модель раздельного кодирования подходит для первичного топ-K отбора без внешних платежей за использование интерфейса провайдера, но требует локального центрального процессора и оперативной памяти и не заменяет полную интерпретацию большой языковой модели.

Расчет внедрения строится от измеримых показателей качества и времени. В качестве технического источника времени используется полный прогон оценки ранжирования на файле тестовой выборки \path|data/splits/test.tsv|, где каждая строка задает пару резюме и вакансии с меткой релевантности: файл \path|experiments/results/full_test_ranking_metrics_stdout.txt| фиксирует 521 обработанную текстовую пару текст вакансии и текст резюме, и время выполнения runtime\_seconds = 3.870512. Отсюда среднее модельное время на одну заявку:

\begin{equation}
\label{eq:t_model_sec-16}
T_{\text{model\_sec}} = R_{\text{test}} / N_{\text{test}} = 3.870512 / 521 = 0.007429 \text{с}
\end{equation}

\begin{equation}
\label{eq:t_model_min-17}
T_{\text{model\_min}} = T_{\text{model\_sec}} / 60 = 0.000124 \text{мин.}
\end{equation}

Для ручного сценария в расчет введены явные допущения, поскольку в логах системы автоматизированного подбора персонала под названием Reqcore отсутствует независимый замер рабочего времени специалиста по подбору персонала. В работе для ручной оценки зафиксируем ориентир 45 минут на цикл из 20 заявок от соискателей. Для расчета принято, что один такой цикл содержит ровно 20 заявок от соискателей. После внедрения модуль ранжирования цифровых профилей кандидатов не принимает кадровое решение автономно: он сортирует заявки и передает результат специалисту по подбору персонала. Для контрольной проверки ранжированного списка принята верхняя граница: 4 минуты на 20 заявок.

\begin{equation}
\label{eq:t_human-18}
T_{\text{human}} = T_{\text{manual\_batch}} / N_{\text{batch}} = 45 / 20 = 2.25 \text{мин.}
\end{equation}

\begin{equation}
\label{eq:t_check_hr-19}
T_{\text{check\_hr}} = T_{\text{check\_batch}} / N_{\text{batch}} = 4 / 20 = 0.20 \text{мин.}
\end{equation}

\Needspace{10\baselineskip}
\begingroup
\setlength{\tabcolsep}{3pt}
\renewcommand{\arraystretch}{1}
{\linespread{1}\selectfont
\begin{longtable}{|P{0.22\textwidth}|P{0.28\textwidth}|P{0.28\textwidth}|P{0.16\textwidth}|}
\caption{Сравнение времени ручной обработки и модельного контура}\label{tbl:deploy-time}\\*
\hline
Показатель & Ручная обработка & Модельный контур & Разница \\
\hline
\endfirsthead
\caption[]{Сравнение времени ручной обработки и модельного контура (продолжение)}\\*
\hline
Показатель & Ручная обработка & Модельный контур & Разница \\
\hline
\endhead
\hline
\endfoot
\hline
\endlastfoot
Время на одну заявку & 2.25 мин & 0.000124 мин вывода модели & 2.249876 мин \\
\hline
Время специалиста по подбору персонала после внедрения & 2.25 мин & 0.20 мин проверки & экономия 2.05 мин \\
\hline
Accuracy относительно меток специалистов по подбору персонала & 1.0000 & 0.9750 & -0.0250 \\
\hline
Средняя точность (Average Precision) на тестовой выборке & не применяется & 0.9921 & модельная метрика \\
\hline
Precision@3 для вакансий с хотя бы одним исторически продвинутым кандидатом & не применяется & 0.9958 & модельная метрика \\
\end{longtable}
}
\endgroup

Количество заявок за расчетный месяц принято равным 521, то есть объему отложенного test split. Такой выбор не переносит тестовый набор в production как факт нагрузки, но позволяет сделать воспроизводимый расчет на тех же данных, на которых измерены качество и время. Экономия времени и денежные показатели задаются следующими формулами:

\begin{equation}
\label{eq:s_time_one-20}
S_{\text{time\_one}} = T_{\text{human}} - T_{\text{model\_min}}
\end{equation}

\begin{equation}
\label{eq:s_time_period-21}
S_{\text{time\_period}} = S_{\text{time\_one}} \cdot N_{\text{period}}
\end{equation}

\begin{equation}
\label{eq:c_manual_one-22}
C_{\text{manual\_one}} = (T_{\text{human}} / 60) \cdot R_{\text{hr}}
\end{equation}

\begin{equation}
\label{eq:c_after_one-23}
C_{\text{after\_one}} = (T_{\text{model\_sec}} / 3600) \cdot R_{\text{cpu}} + (T_{\text{check\_hr}} / 60) \cdot R_{\text{hr}}
\end{equation}

\begin{equation}
\label{eq:s_cost_one-24}
S_{\text{cost\_one}} = C_{\text{manual\_one}} - C_{\text{after\_one}}
\end{equation}

\begin{equation}
\label{eq:s_cost_period-25}
S_{\text{cost\_period}} = S_{\text{cost\_one}} \cdot N_{\text{period}}
\end{equation}

\begin{equation}
\label{eq:c_support_month-26}
C_{\text{support\_month}} = H_{\text{support}} \cdot R_{\text{engineer}}
\end{equation}

\begin{equation}
\label{eq:c_impl-27}
C_{\text{impl}} = H_{\text{impl}} \cdot R_{\text{engineer}} + (T_{\text{train\_sec}} / 3600) \cdot R_{\text{cpu}}
\end{equation}

\begin{equation}
\label{eq:e_month-28}
E_{\text{month}} = S_{\text{cost\_period}} - C_{\text{support\_month}}
\end{equation}

\begin{equation}
\label{eq:p_payback-29}
P_{\text{payback}} = C_{\text{impl}} / E_{\text{month}}
\end{equation}

\begin{equation}
\label{eq:roi_12-30}
ROI_{12} = ((E_{\text{month}} \cdot 12) - C_{\text{impl}}) / C_{\text{impl}} \cdot 100 \%
\end{equation}

где S\_time\_one - экономия времени на одной заявке, S\_time\_period - экономия времени за период, N\_period - число заявок за расчетный период, C\_manual\_one - стоимость ручной обработки одной заявки, C\_after\_one - стоимость обработки одной заявки после внедрения, S\_cost\_one - денежная экономия на одной заявке, S\_cost\_period - денежная экономия за период, R\_hr - часовая ставка специалиста по подбору персонала, R\_cpu - стоимость часа сервера с центральным процессором, R\_engineer - инженерная ставка, H\_support - часы ежемесячного сопровождения, H\_impl - часы разового внедрения, T\_train\_sec - время первичного обучения модели в секундах, C\_support\_month - ежемесячные затраты на сопровождение, C\_impl - затраты на внедрение, E\_month - чистый ежемесячный эффект, P\_payback - срок окупаемости, ROI\_12 - возврат инвестиций за 12 месяцев.

Подстановка значений для времени:

\begin{equation}
\label{eq:s_time_one-31}
S_{\text{time\_one}} = 2.25 - 0.000124 = 2.249876 \text{мин.}
\end{equation}

\begin{equation}
\label{eq:s_time_period-32}
S_{\text{time\_period}} = 2.249876 \cdot 521 = 1172.19 \text{мин}= 19.54 \text{ч за расчетный месяц}
\end{equation}

Для денежного расчета заданы следующие прозрачные допущения: часовая ставка специалиста по подбору персонала 600 руб./ч; стоимость машинного времени сервера с центральным процессором 30 руб./ч; инженерная ставка для внедрения и сопровождения 1500 руб./ч системы; разовое внедрение готового прототипа в пилотный контур системы автоматизированного подбора персонала под названием Reqcore занимает 16 инженерных часов; ежемесячное сопровождение занимает 2 инженерных часа. Стоимость первоначального обучения модели также включена в затраты на внедрение: внешний замер команды обучения равен 6254.096 с, что при ставке 30 руб./ч дает 52.12 руб.

Расчетные денежные показатели:

\Needspace{10\baselineskip}
\begingroup
\setlength{\tabcolsep}{3pt}
\renewcommand{\arraystretch}{1}
{\linespread{1}\selectfont
\begin{longtable}{|P{0.30\textwidth}|P{0.34\textwidth}|P{0.28\textwidth}|}
\caption{Расчётные денежные показатели внедрения}\label{tbl:roi}\\*
\hline
Метрика & Расчет & Значение \\
\hline
\endfirsthead
\caption[]{Расчётные денежные показатели внедрения (продолжение)}\\*
\hline
Метрика & Расчет & Значение \\
\hline
\endhead
\hline
\endfoot
\hline
\endlastfoot
Стоимость ручной обработки одной заявки & \texttt{(2.25 / 60) * 600} & 22.50 руб. \\
\hline
Стоимость после внедрения на одну заявку & \texttt{(0.007429 / 3600) * 30 + (0.20 / 60) * 600} & 2.0001 руб. \\
\hline
Экономия на одной заявке & 22.50 - 2.0001 & 20.4999 руб. \\
\hline
Экономия за расчетный месяц & \texttt{20.4999 * 521} & 10680.47 руб. \\
\hline
Затраты на сопровождение за месяц & \texttt{2 * 1500} & 3000.00 руб. \\
\hline
Чистый ежемесячный эффект & 10680.47 - 3000.00 & 7680.47 руб. \\
\hline
Затраты на внедрение & \texttt{16 * 1500 + (6254.096 / 3600) * 30} & 24052.12 руб. \\
\hline
Срок окупаемости & \texttt{24052.12 / 7680.47} & 3.13 месяца \\
\hline
Чистый эффект за 12 месяцев & \texttt{7680.47 * 12} & 92165.61 руб. \\
\hline
ROI за 12 месяцев & (92165.61 - 24052.12) / 24052.12 * 100\% & 283.2\% \\
\end{longtable}
}
\endgroup

Практическая ценность внедрения состоит в том, что модель сокращает первичный машинно-ручной цикл обработки с 2.25 минуты ручной работы до 0.20 минуты проверки специалистом по подбору персонала и 0.007429 секунды локального вывода модели на заявку. При сохранении контрольной роли специалиста по подбору персонала качество модели на отложенной тестовой выборке остается высоким: доля правильных решений (Accuracy) равна 0.9750, средняя точность по кривой точность-полнота (Average Precision) равна 0.9921, а точность на первых трех позициях (Precision@3) равна 0.9958. Следовательно, экономический эффект достигается не за счет исключения человека из процесса, а за счет переноса трудоемкого первичного просмотра на локальный модуль ранжирования цифровых профилей соискателей с последующей быстрой проверкой результата специалистом по подбору персонала.

\section{Валидация трех стратегий ранжирования на отложенных данных}

В данной работе вместо продуктового сравнительного A/B-эксперимента используется оценка качества трех стратегий ранжирования на фиксированном отложенном наборе данных. Такой выбор соответствует текущей архитектуре системы автоматизированного подбора персонала под названием Reqcore: конфигурация искусственного интеллекта хранится на уровне организации в таблице ai\_config, где задаются активный провайдер, модель и ценовые параметры отправки и получения знаков, а маршрут POST \path|/api/applications/[id]/analyze| при каждом запуске анализа читает одного активного провайдера provider и одну активную модель для всей выбранной организации. Следовательно, в текущей реализации без дополнительного серверного механизма назначения экспериментальной группы невозможно корректно провести одновременный A/B-тест внутри одной и той же организации на уровне отдельных вакансий или заявок. Поэтому наиболее корректным и простым способом показать эффективность является сравнение на отложенных данных вероятностной функции ранжирования документов под названием BM25, перекрестной модели совместного кодирования текстовой пары (cross-encoder) и модели раздельного кодирования (bi-encoder) на заранее сформированных файлах обучающей, валидационной и тестовой выборок: train.tsv, val.tsv и test.tsv.

В качестве основного сценария оценки используется сравнение ровно трех моделей. Вероятностная функция ранжирования документов под названием BM25 не обучается и проверяет силу лексического совпадения терминов, навыков и должностей. Перекрестная модель совместного кодирования текстовой пары под названием \path|cointegrated/rubert-tiny2| (cross-encoder) обучается как модель классификации последовательностей для предсказания класса кадрового решения по паре и возвращает вероятность label = 1 для входа текстов [vacancy\_text, resume\_text]. Модель раздельного кодирования (bi-encoder), сохраненная в каталоге \path|experiments/models/bi_encoder_rubert_tiny2|, преобразует вакансию и резюме в отдельные плотные векторные представления, после чего вычисляет косинусное сходство (cosine similarity); именно эта логика соответствует текущей ветке, где выполняется условие provider === 'biencoder', и числовой оценке, записываемой в поле \path|application.score| в системе автоматизированного подбора персонала под названием Reqcore.

Преимущество такого подхода состоит в том, что он не требует нового эксперимента в реальном времени на пользователях и не требует обязательного привлечения дополнительных экспертных разметчиков. В качестве источника истинности используются уже подготовленные исторические текстовые пары текст резюме и текст вакансии с метками, полученными из исходных событий от сотрудников по подбору персонала и затем разделенные на выборки train, val и test. Это позволяет проверить не только абсолютное качество моделей, но и рациональность компромисса между качеством, скоростью, стоимостью вывода модели и пригодностью для системы автоматизированного подбора персонала.

Офлайн-оценку целесообразно разделить на два уровня.

\Needspace{10\baselineskip}
\begingroup
\setlength{\tabcolsep}{3pt}
\renewcommand{\arraystretch}{1}
{\linespread{1}\selectfont
\begin{longtable}{|P{0.30\textwidth}|P{0.34\textwidth}|P{0.28\textwidth}|}
\caption{Уровни офлайн-валидации трёх стратегий ранжирования}\label{tbl:offline}\\*
\hline
Уровень оценки & Что измеряется & Артефакты \\
\hline
\endfirsthead
\caption[]{Уровни офлайн-валидации трёх стратегий ранжирования (продолжение)}\\*
\hline
Уровень оценки & Что измеряется & Артефакты \\
\hline
\endhead
\hline
\endfoot
\hline
\endlastfoot
Бинарная классификация пар (pair-level classification) & Качество отделения пар, где кандидат был принят к рассмотрению или приглашен на следующий этап, от пар, где был зафиксирован отказ, для перекрестной модели совместного кодирования и модели раздельного кодирования (bi-encoder) & train.tsv, val.tsv, test.tsv \\
\hline
Ранжирование внутри вакансии (vacancy-level ranking) & Качество сортировки кандидатов внутри одной вакансии для вероятностной функции ранжирования документов под названием BM25, перекрестной модели совместного кодирования и модели раздельного кодирования (bi-encoder) & группировка test.tsv по vacancy\_id \\
\end{longtable}
}
\endgroup

На уровне отдельных используются метрики бинарного разделения: средняя точность (Average Precision), гармоническая средняя (F1-score), доля правильных ответов (Accuracy) и коэффициент корреляции Мэтьюса (Matthews Correlation Coefficient, MCC). Они количественно отражают способность вероятностной оценки перекрестной модели совместного кодирования и косинусного сходства (cosine similarity) модели раздельного кодирования различать текстовые пары текст вакансии и текст резюме, по которым специалист по подбору персонала продвинул кандидата, и пары, по которым был зафиксирован отказ. На уровне ранжирования внутри вакансии используются точность на первых K позициях (Precision at K, Precision@K), полнота на первых K позициях (Recall at K, Recall@K), средняя обратная позиция первого исторически продвинутого кандидата (Mean Reciprocal Rank, MRR) и нормализованная дисконтированная кумулятивная выгода на первых K позициях (Normalized Discounted Cumulative Gain at K, NDCG@K) после группировки по идентификатору вакансии vacancy\_id и сортировки кандидатов по числовой оценке для каждой модели, которую мы хотим сравнить. Именно этот уровень оценки наиболее близок к реальному использованию модуля как инструмента предварительного отбора и сортировки цифровых профилей кандидатов.

Минимальная процедура эксперимента выглядит следующим образом: обучаемые модели используют файл train.tsv, пороговые и вспомогательные параметры подбираются на валидационной выборке val.tsv, а финальные метрики публикуются только на тестовой выборке test.tsv, которая не используется ни для обучения, ни для подбора решений. Для оценки ранжирования каждая вакансия рассматривается как отдельный запрос, а все связанные с ней резюме сортируются по числовой оценке вероятностной функции ранжирования документов под названием BM25, вероятности перекрестной модели совместного кодирования или косинусному сходству (cosine similarity) модели раздельного кодирования. Такой дизайн позволяет получить воспроизводимую и технически корректную оценку качества без вмешательства в рабочий продуктовый контур системы автоматизированного подбора персонала под названием Reqcore.

Ограничения данного подхода также должны быть зафиксированы явно. Во-первых, это валидация на отложенных размеченных данных, а не эксперимент в реальном времени влияния на бизнес-метрики подбора персонала. Во-вторых, используемые метки основаны на исторических событиях из системы по подбору персонала и потому являются приближением к реальным закономерностям предметной области подбора персонала, а не независимой экспертной эталонной разметкой (gold labels). В-третьих, экспериментальная ветка обработки biencoder в текущем коде возвращает только числовую оценку и не формирует объяснения на уровне отдельных критериев, включая поля evidence, strengths и gaps, поэтому оценка на отложенных данных должна интерпретироваться именно как проверка качества числового ранжирования, а не как сравнение полноты объяснимости с веткой, где мы используем одного из провайдеров большой языковой модели для оценки цифрового профиля кандидата.

Если в дальнейшем потребуется дополнительная качественная проверка, ее можно выполнить на небольшой подвыборке вакансий в формате экспертной проверки первых K кандидатов (top-K review). Однако для базовой демонстрации эффективности в данной работе достаточно именно оценки на фиксированной тестовой выборке (test split), поскольку она лучше соответствует текущей архитектуре продукта и не требует вводить искусственный A/B-дизайн там, где он пока не поддержан на уровне серверной логики.

\section{Выводы по четвертому разделу}

\begin{compactenum}
  \item Экспериментальная оценка показывает, что модель раздельного кодирования является масштабируемым компромиссом между качеством ранжирования и стоимостью вывода по сравнению с более дорогими попарными стратегиями.
  \item Интеграционная проверка подтверждает, что модуль можно подключить к серверному контуру Reqcore через HTTP-интерфейс и использовать для записи итогового числового балла заявки.
  \item Экономический расчет показывает потенциальное сокращение трудоемкости первичного просмотра кандидатов при сохранении контрольной роли специалиста по подбору персонала.
  \item Офлайн-валидация трех стратегий ранжирования фиксирует текущие границы доказанности результата: качество подтверждено на отложенных данных, но не заменяет промышленный A/B-эксперимент и аудит справедливости.
\end{compactenum}
